\chapter{Gaussian Projection and High-Order Quadrature Filters}
\label{ch:bf-highdim-gaussian-quadrature}

\section{Chapter Claim}
\label{sec:bf-highdim-gaussian-claim}

Gaussian projection filters remain the most inspectable first rung for
high-dimensional nonlinear SSMs, but unstructured high-order quadrature is not
scalable.  The defensible industrial direction is block-local or low-rank
quadrature with explicit approximation diagnostics, not global CUT4 on the full
state.

\section{Assumptions}
\label{sec:bf-highdim-gaussian-assumptions}

This chapter assumes that Gaussian projection is being used as an approximation
to a nonlinear filtering law, not as an exact statement unless the model is
affine Gaussian or another exactness argument is supplied.  Sparse-grid and
block-local discussion assumes exploitable locality, sparsity, or low effective
interaction order; without such structure these methods remain candidate
approximations with no promotion claim.

\section{Moment Projection}
\label{sec:bf-highdim-moment-projection}

Assume the predictive state is approximated by
$x_t\mid y_{1:t-1}\approx \mathcal N(m_t^-,P_t^-)$.  A nonlinear observation
map $h_\theta$ is reduced to approximate moments
\begin{align}
  \hat y_t &= \mathbb E[h_\theta(X_t)],\\
  S_t &= \operatorname{Var}(h_\theta(X_t)) + R_\theta,\\
  C_t &= \operatorname{Cov}(X_t,h_\theta(X_t)).
\end{align}
The Gaussian update is
\begin{align}
  K_t &= C_t S_t^{-1},\\
  m_t &= m_t^- + K_t(y_t-\hat y_t),\\
  P_t &= P_t^- - K_t S_t K_t^\top.
\end{align}
The approximation is entirely in the moment calculation.  The update is the
Gaussian projection update for those moments.

\section{Derivative-Based Filters}
\label{sec:bf-highdim-derivative-filters}

The extended Kalman filter uses first derivatives of the transition and
observation maps.  An iterated EKF repeats local linearization around an updated
state estimate and is often a better optimizer or local filter when the
observation equation is strongly nonlinear.  A second-order EKF can include
Hessian corrections to the mean and covariance.  These filters scale better in
state dimension than dense high-order quadrature when Jacobian and Hessian
structure is sparse or block-local, but their evidence burden shifts to
derivative correctness and local approximation error.

\section{Sigma-Point Rules}
\label{sec:bf-highdim-sigma-rules}

A sigma-point rule replaces expectations by
\[
  \mathcal Q[\varphi] = \sum_{j=1}^{N_q} w_j \varphi(m + L\xi_j),
  \qquad LL^\top=P.
\]
BayesFilter currently exposes TensorFlow SVD/eigen sigma-point backends for
cubature, unscented, and CUT4-G value paths.  Chapter~\ref{ch:bf-sigma-point-filters}
records the CUT4-G rule:
\[
  N_{\mathrm{CUT4}}(d)=2d+2^d.
\]
This is the central high-dimensional warning.  At augmented dimension $d=15$,
the rule needs $32798$ points; at $d=100$, the corner set is physically
irrelevant as an implementation proposal.

\begin{remark}[High-order does not mean high-dimensional]
\label{rem:bf-highdim-cut4-boundary}
High polynomial exactness at fixed dimension is not the same as high-dimensional
tractability.  CUT4-G is useful as a small-dimensional reference and a block
rule.  It is not an unstructured industrial default.
\end{remark}

\section{Sparse Grids And Block Rules}
\label{sec:bf-highdim-sparse-block}

Sparse-grid and Smolyak rules try to spend points on lower-order interactions
instead of all tensor-product corners.  Block-local quadrature goes further:
partition the state or innovation coordinates into groups, apply higher-order
rules inside each group, and stitch the result through a Gaussian or ensemble
approximation.  This changes the approximation class.  It requires diagnostics
for neglected cross-block interactions, covariance positive-definiteness, and
downstream likelihood or posterior distortion.

\section{Algorithmic Template}
\label{sec:bf-highdim-block-quadrature-template}

\begin{enumerate}[label=(\arabic*)]
  \item Choose a stochastic integration block partition.
  \item For each block, select EKF, CKF/UKF, CUT4, or sparse-grid moment rule.
  \item Propagate block moments through the nonlinear transition/observation
    maps.
  \item Reconstruct global mean/covariance with an explicit cross-block policy.
  \item Record point count, PSD residual, cross-block residual proxy, value
    parity on small cells, and runtime.
\end{enumerate}

\section{BayesFilter Evidence Box}
\label{sec:bf-highdim-gaussian-evidence}

The existing V1 record supports finite value and score diagnostics for Models
A-C on narrow cells and CPU/GPU/XLA timing diagnostics for exact tested Model B
rows.  The P8 harness in this program is allowed to add point-count and
execution diagnostics for block extensions.  Neither artifact validates
high-dimensional posterior accuracy.

\section{Source And Evidence Ledger}
\label{sec:bf-highdim-gaussian-ledger}

\begin{tabular}{p{0.28\linewidth}p{0.24\linewidth}p{0.38\linewidth}}
\toprule
Claim & Support class & Boundary \\
\midrule
CUT4-G point count is $2d+2^d$ & BayesFilter implementation and existing
Chapter~\ref{ch:bf-sigma-point-filters} derivation & Applies to implemented
CUT4-G rule, not all high-order rules. \\
Gaussian projection update formulas & standard derivation, chapter exposition &
Moment approximation quality is not guaranteed. \\
Sparse/block quadrature is a candidate direction & literature taxonomy and
research synthesis & No BayesFilter default or production claim. \\
\bottomrule
\end{tabular}

\section{Unresolved Claim Register}
\label{sec:bf-highdim-gaussian-unresolved}

\begin{itemize}
  \item No sparse-grid backend is implemented in this phase.
  \item No block-local quadrature theorem is proved for NAWM.
  \item No default policy change is justified.
\end{itemize}

\section{What Is Not Concluded}
\label{sec:bf-highdim-gaussian-nonclaims}

The chapter does not conclude that EKF, sparse grids, CUT4, or block quadrature
is accurate for a given high-dimensional DSGE model.  It ranks them as
diagnosable candidates and defines the evidence needed before promotion.
